\section{Introducción}
La presente práctica de Física General tiene como propósito el estudio y análisis del movimiento 
de los cuerpos a través de los conceptos fundamentales de la cinemática. En este contexto, 
se abordan magnitudes físicas como la posición, el desplazamiento, la velocidad y la aceleración, 
las cuales permiten describir el comportamiento de un objeto en función del tiempo 
sin considerar las causas que originan dicho movimiento.\\
\vspace{10pt}

Asimismo, se desarrollan y aplican los modelos matemáticos correspondientes al Movimiento Rectilíneo Uniforme (MRU), 
al Movimiento Rectilíneo Uniformemente Acelerado (MRUA) 
y a la caída libre, estableciendo relaciones entre las variables físicas involucradas. 
El uso del cálculo diferencial e integral permite una comprensión más profunda de estos fenómenos, 
facilitando la obtención de ecuaciones que describen el movimiento de manera rigurosa.\\
\vspace{10pt}

Finalmente, se complementa el análisis teórico mediante el uso de herramientas tecnológicas, 
como simuladores, que permiten visualizar e interpretar gráficamente los resultados obtenidos, 
fortaleciendo así la comprensión conceptual y la validación de los procedimientos desarrollados.